\section{实验阅读法：从现象回到假设}
这类 20 分钟短课最大的价值，往往不是多讲了一个公式，而是把一个实验现象钉住。真正复看时，应该问三个问题：实验在验证什么假设、用了什么代理指标、如果结论不成立会怎样。这样才能避免把一次性的实验结果误读为普适规律。

\begin{knowledgebox}{读实验结果的三个抓手}
\begin{itemize}
  \item 先看自变量：究竟改变了什么配置
  \item 再看观测量：reward、成功率还是 loss
  \item 最后看外推边界：这个结论能否推广到别的环境
\end{itemize}
\end{knowledgebox}

\subsection{本章小结}
实验分析的重点不只是“结果是什么”，而是“这个结果支持了哪一个训练判断”。这是后续做 Agentic RL ablation 时最需要保留的习惯。

\newpage
